\section{Прогноз и итоговый вывод}

\subsection{Прогноз для заданного значения \(x^*\)}

Заданное прогнозное значение фактора: \(x^* = 1207.9904\) МБ/с (что составляет 107\% от среднего уровня скорости чтения \(\bar{x} = 1128.963\) МБ/с).

В таблице~\ref{tab:prediction} приведены прогнозные значения \(\hat{y}(x^*)\) по каждой из трёх моделей.

\begin{table}[H]
\centering
\caption{Прогнозные значения \(\hat{y}(x^*)\) для \(x^* = 1207.9904\) МБ/с}
\label{tab:prediction}
\begin{tabular}{lc}
\toprule
Модель & \(\hat{y}(x^*)\), МБ/с \\
\midrule
Линейная \(y = a + bx\) & 1169.04 \\
Квадратичная \(y = a + bx + cx^2\) & 1290.63 \\
Степенная \(y = ax^b\) & 1116.07 \\
\bottomrule
\end{tabular}
\end{table}

\subsection{Выбор модели для прогноза}

В качестве основной модели для прогноза выбирается \textbf{степенная модель} \(\hat{y} = 0.6945 \, x^{1.0402}\). Обоснование выбора:

\begin{enumerate}
    \item степенная модель имеет наименьшую среднюю ошибку аппроксимации (\(A = 16.74\%\));
    \item остатки степенной модели наиболее близки к случайным (25 смен знака, наибольшее среди трёх моделей), а в области малых \(x\) средний остаток близок к нулю (\(-0.75\));
    \item показатель степени \(\hat{b} = 1.0402 \approx 1\) указывает на почти пропорциональную зависимость, что является естественным и содержательно интерпретируемым для физических характеристик SSD;
    \item квадратичная модель, имеющая чуть более высокий \(R^2\), содержит отрицательный квадратичный член, приводящий к убыванию прогноза при больших \(x\), что противоречит физическому смыслу: скорость записи SSD не должна снижаться с ростом скорости чтения.
\end{enumerate}

Таким образом, прогнозное значение скорости записи SSD при скорости чтения \(x^* = 1207.9904\) МБ/с (на 7\% выше среднего уровня) составляет:

\[
\hat{y}(x^*) = 0.6945 \cdot (1207.9904)^{1.0402} \approx 1116.07 \text{ МБ/с}.
\]

\subsection{Итоговый вывод}

В рамках РГР №3 (вариант A-5) исследовалась зависимость скорости записи SSD от скорости чтения на основе реального набора данных \texttt{openintro::ssd\_speed} (\(n = 54\), \(\alpha = 0.05\)).

По диаграмме рассеяния была выдвинута гипотеза о положительной связи, близкой к линейной. Построены три регрессионные модели: линейная \(\hat{y} = -24.358 + 0.9879x\), квадратичная \(\hat{y} = -291.300 + 1.6032x - 0.000243x^2\) и степенная \(\hat{y} = 0.6945x^{1.0402}\). Все три модели показали умеренное качество аппроксимации (\(R^2 \approx 0.74\text{--}0.75\), \(A \approx 17\text{--}19\%\)).

По совокупности критериев (\(R^2\), средней ошибке аппроксимации и анализу остатков) \textbf{степенная модель признана предпочтительной}: она имеет наименьшую среднюю ошибку аппроксимации (\(A = 16.74\%\)), наиболее равномерные остатки и физически осмысленное поведение во всём диапазоне \(x\). Линейная модель уступает по качеству приближения, но статистически значима (\(t_{\text{набл}} = 12.31\), \(p < 0.0001\)), что подтверждает наличие положительной связи.

Для \(x^* = 1207.9904\) МБ/с (107\% от среднего) прогноз скорости записи по степенной модели составляет \(\hat{y} \approx 1116\) МБ/с. Содержательно это означает, что скорость записи растёт почти пропорционально скорости чтения с коэффициентом пропорциональности около 0.7 при средних значениях и с небольшим опережением при высоких. Результаты могут использоваться для прогнозирования характеристик SSD с учётом выявленной гетероскедастичности при высоких скоростях.
